\documentclass[10pt,a4paper]{article}

%----------------------------------------------------------------------------------------
%	PACKAGES
%----------------------------------------------------------------------------------------
\usepackage[left=0.7in,top=0.7in,right=0.7in,bottom=0.7in]{geometry}
\usepackage[utf8]{inputenc}
\usepackage[T1]{fontenc}
\usepackage[default]{sourcesanspro}
\usepackage{xcolor}
\usepackage{titlesec}
\usepackage{enumitem}
\usepackage{parskip}
\usepackage{hyperref}
\usepackage{fontawesome5}

%----------------------------------------------------------------------------------------
%	COLORS & STYLES
%----------------------------------------------------------------------------------------
\definecolor{titleblue}{HTML}{00199e}
\definecolor{subtitleblue}{HTML}{2ec1e0}
\definecolor{darktext}{HTML}{222222}
\definecolor{linkblue}{HTML}{0B5FA5}
\definecolor{subtle}{HTML}{666666}

\color{darktext}
\setlength{\parskip}{0.1em}

\titleformat{\section}{\normalsize\bfseries\color{titleblue}}{}{0em}{}[\vspace{-0.3em}\rule{\textwidth}{0.4pt}]
\titlespacing*{\section}{0pt}{0.6em}{0.2em}

\newcommand{\blueitem}[1]{\textcolor{subtitleblue}{\textbf{#1}}}
\newcommand{\sub}[1]{\textcolor{subtle}{\footnotesize #1}}

\setlist[itemize]{label=\textbullet, leftmargin=*, noitemsep, topsep=2pt, parsep=0pt}

\hypersetup{colorlinks=true, urlcolor=linkblue, linkcolor=linkblue}
%----------------------------------------------------------------------------------------
%	DOCUMENT CONTENT
%----------------------------------------------------------------------------------------
\begin{document}

%--- HEADER ---
{\Large \textbf{\textcolor{titleblue}{Iván Sevillano-García}}} \hfill {\footnotesize \textcolor{subtle}{521 citas \textbar\ índice h 4 \textbar\ 512 desde 2021}} \\[0.1em]
{\footnotesize Instituto DaSCI, Universidad de Granada, España} \\[0.1em]
{\footnotesize
\faEnvelope\ \href{mailto:isevillano@ugr.es}{isevillano@ugr.es} \quad
\faGlobe\ \href{https://www.linkedin.com/in/iván-sevillano-garcía-304801116}{LinkedIn} \quad
\faGithub\ \href{https://github.com/isega24}{GitHub} \quad
\faGraduationCap\ \href{https://scholar.google.com/citations?user=g5xyauEAAAAJ}{Google Scholar} \quad
\faOrcid\ \href{https://orcid.org/0000-0002-5029-9106}{ORCID}
}

\vspace{0.3em}

%--- PROFILE ---
\section*{Perfil}

Investigador posdoctoral en IA Confiable --- IA explicable (XAI), detección de datos fuera de distribución, evaluación y auditoría de explicaciones, y mejora de modelos guiada por explicaciones. Doctor Cum Laude (UGR, 2025); publicaciones en \textit{Pattern Recognition}, \textit{Machine Intelligence Research}, \textit{Information Fusion} e \textit{IEEE JBHI}. Mis intereses de investigación actuales incluyen IA agéntica, ingeniería de \textit{harness} e ingeniería de bucles (\textit{loop engineering}).

%--- EDUCATION ---
\section*{Formación}

\blueitem{Doctorado en Tecnologías de la Información y la Comunicación \hfill 2021--2025} \\
\sub{Universidad de Granada --- Instituto DaSCI} \\
\textit{Tesis:} Explainable AI: Methods for the Generation, Evaluation, and Correction of Explanations \\
\sub{Directores: F.~Herrera, J.~Luengo \quad Defensa: 13 nov 2025 \quad \textbf{Cum Laude}} \\
\sub{Contribuciones: REVEL (evaluación de explicaciones), X-SHIELD (regularización de explicaciones), STOOD-X (detección OOD explicable).} \\
\sub{\href{https://hdl.handle.net/10481/112935}{hdl.handle.net/10481/112935} \quad ISBN 9791370150433}

\vspace{0.2em}

\blueitem{Máster en Ciencia de Datos e Ingeniería de Computadores \hfill 2019--2020} \\
\sub{Universidad de Granada}

\vspace{0.2em}

\blueitem{Grado en Matemáticas \hfill 2013--2019} \\
\sub{Universidad de Granada --- ETSIIT}

\vspace{0.2em}

\blueitem{Grado en Ingeniería Informática \hfill 2013--2019} \\
\sub{Universidad de Granada --- ETSIIT}

%--- RESEARCH EXPERIENCE ---
\section*{Experiencia investigadora}

\blueitem{Investigador predoctoral \hfill 2019--2025} \\
\sub{Universidad de Granada --- Instituto DaSCI}
\begin{itemize}
    \item \textbf{REVEL:} diseño de un marco teórico para auditar explicaciones lineales locales en robustez, concisión y prescriptividad; evaluado sobre \textit{benchmarks} de clasificación de imágenes con aprendizaje profundo.
    \item \textbf{X-SHIELD:} propuesta de una familia de regularizadores guiados por explicaciones que seleccionan características de entrada que ocultar durante el entrenamiento, mejorando simultáneamente rendimiento y explicabilidad del modelo.
    \item \textbf{STOOD-X:} desarrollo de un método de detección OOD explicable en dos etapas que combina el test no paramétrico de Wilcoxon-Mann-Whitney con explicaciones visuales basadas en conceptos.
    \item \textbf{COVIDGR:} codiseño de un conjunto de radiografías de tórax COVID-19 equilibrado en todos los niveles de gravedad y de la metodología COVID-SDNet.
\end{itemize}

\blueitem{Investigador visitante \hfill feb--mar 2026} \\
\sub{Universidad Politécnica de Poznań, Polonia --- anfitrión: Prof.~J.~Stefanowski}
\begin{itemize}
    \item Investigación sobre explicaciones basadas en prototipos y contrafactuales para métodos XAI post-hoc.
\end{itemize}

\blueitem{Colaboración industrial --- ArcelorMittal \hfill 2020--2021} \\
\sub{vía DaSCI, UGR}
\begin{itemize}
    \item Codesarrollo de un \textit{ensemble} de aprendizaje profundo para la segmentación de imágenes metalográficas; publicación conjunta del conjunto de datos MetalDAM.
\end{itemize}

%--- PUBLICATIONS ---
\section*{Publicaciones}

\small

\blueitem{Revistas}

\textbf{STOOD-X: Explainable OOD Detection via Nonparametric Statistical Testing on Large-Scale Datasets.} \\
\sub{\textit{Pattern Recognition}, 2026. \textbf{I.~Sevillano-García}, J.~Luengo, F.~Herrera. \href{https://doi.org/10.1016/j.patcog.2026.113254}{DOI} \quad (3 citas)}

\textbf{X-SHIELD: Regularization for Explainable Artificial Intelligence.} \\
\sub{\textit{Machine Intelligence Research}, 23, 528--544, 2026. \textbf{I.~Sevillano-García}, J.~Luengo, F.~Herrera. \href{https://doi.org/10.1007/s11633-025-1576-y}{DOI} \quad (9 citas)}

\textbf{A Three-Level Framework for LLM-Enhanced eXplainable AI.} \\
\sub{\textit{Information Systems Frontiers}, 2025. M.~Bello, R.~Bello, M.M.~García, A.~Nowé, \textbf{I.~Sevillano-García}, F.~Herrera. \href{https://doi.org/10.1007/s10796-025-10668-1}{DOI} \quad (14 citas)}

\textbf{REVEL Framework to Measure Local Linear Explanations for Black-Box Models.} \\
\sub{\textit{Intl.~J.~Intelligent Systems}, 2023. \textbf{I.~Sevillano-García}, J.~Luengo, F.~Herrera. \href{https://doi.org/10.1155/2023/8068569}{DOI} \quad (15 citas)}

\textbf{A Tutorial on Segmentation of Metallographic Images: MetalDAM Dataset, DL Ensemble.} \\
\sub{\textit{Information Fusion}, 2022. J.~Luengo, R.~Moreno, \textbf{I.~Sevillano-García}, et~al. \href{https://doi.org/10.1016/j.inffus.2021.09.018}{DOI} \quad (59 citas)}

\textbf{COVIDGR Dataset \& COVID-SDNet Methodology for Predicting COVID-19 (Chest X-ray).} \\
\sub{\textit{IEEE JBHI}, 24(12), 3595--3605, 2020. S.~Tabik, A.~Gómez-Ríos, J.L.~Martín-Rodríguez, \textbf{I.~Sevillano-García}, et~al. \href{https://doi.org/10.1109/JBHI.2020.3037127}{DOI} \quad (477 citas)}

\blueitem{Congresos}

\textbf{Fine-Grained Prostate Cancer Tissue Classification via Expert-Annotated Dataset \& Interpretable DL.} \\
\sub{\textit{IEEE CAI 2026}, 1264--1269. J.~Ayensa-Jiménez, D.~Navarro, \textbf{I.~Sevillano-García}, et~al. \href{https://doi.org/10.1109/CAI68641.2026.11536250}{DOI}}

\textbf{Optimizing LIME Explanations Using REVEL Metrics.} \\
\sub{\textit{HAIS}, 304--313, 2023. \textbf{I.~Sevillano-Garcia}, J.~Luengo, F.~Herrera. \href{https://doi.org/10.1007/978-3-031-40725-3_26}{DOI}}

\textbf{Low-Impact Feature Reduction Regularization Term: How to Improve AI with Explainability.} \\
\sub{\textit{xAI 2023 (Late-breaking Work, poster)}, 135--139. \textbf{I.~Sevillano-García}, J.~Luengo, F.~Herrera. \href{https://dblp.org/rec/conf/xai/Sevillano-Garcia23}{DBLP}}

\textbf{Slicer: Feature Learning for Class Separability with LS-SVM Loss (COVID-19 X-ray).} \\
\sub{\textit{HAIS}, 305--315, 2021. D.~Charte, \textbf{I.~Sevillano-García}, et~al. \href{https://doi.org/10.978-3-030-86271-8_26}{DOI}}

\normalsize

%--- OPEN SOURCE ---
\section*{Software y datos de código abierto}

\blueitem{Software}
\begin{itemize}
    \item \textbf{S-ReVel} --- audita explicaciones lineales locales (p.\,ej.\ LIME) en robustez, concisión y prescriptividad. \href{https://github.com/ari-dasci/S-ReVel}{github.com/ari-dasci/S-ReVel}. AGPL-3.0.
    \item \textbf{S-STOOD-X} --- detección OOD explicable mediante el test de Wilcoxon-Mann-Whitney y explicaciones visuales basadas en conceptos. \href{https://github.com/ari-dasci/S-STOOD-X}{github.com/ari-dasci/S-STOOD-X}. GPL-3.0.
    \item \textbf{S-T-SHIELD} --- regularización guiada por explicaciones que oculta características de entrada durante el entrenamiento para mejorar rendimiento y explicabilidad. \href{https://github.com/ari-dasci/S-T-SHIELD}{github.com/ari-dasci/S-T-SHIELD}. GPL-3.0.
\end{itemize}

\blueitem{Conjuntos de datos}
\begin{itemize}
    \item \textbf{COVIDGR} --- conjunto de radiografías de tórax COVID-19 con clasificación fina por niveles de gravedad, desde RT-PCR positiva sin síntomas hasta casos graves. \href{https://github.com/ari-dasci/OD-covidgr}{github.com/ari-dasci/OD-covidgr}.
    \item \textbf{MetalDAM} --- conjunto de imágenes metalográficas para segmentación microestructural en el control de calidad de la industria del acero. \href{https://github.com/ari-dasci/OD-MetalDAM}{github.com/ari-dasci/OD-MetalDAM}.
\end{itemize}

%--- TEACHING & MENTORING ---
\section*{Docencia y tutorización}

\begin{itemize}
    \item Codirección de 2 tesis doctorales (programa UGR), en curso: temas como \textit{Machine Unlearning}, métricas XAI e IA agéntica.
    \item Codirección de trabajos de fin de grado y de máster, UGR.
    \item Docencia universitaria: Metodología de la Programación (1.er curso de grado, Dpto.~de Ciencias de la Computación e IA, UGR, 2024--2025). \href{https://grados.ugr.es/ramas/ingenieria-arquitectura/grado-ingenieria-informatica/metodologia-la-programacion}{Página de la asignatura}.
\end{itemize}

%--- ACADEMIC SERVICE ---
\section*{Actividad académica}

\blueitem{Ponencias invitadas y docencia:}
\begin{itemize}
    \item Escuela Internacional de Verano ADIA Lab sobre IA Explicable --- ``Lab on xAI Tools'' (Granada, mayo 2025).
    \item Instructor, Seminarios INCIBE sobre IA y Ciberseguridad (DaSCI / Cámara de Comercio, 2025). Coordiné el Bloque 1 sobre IA Confiable junto a F.~Herrera.
\end{itemize}

\blueitem{Revisión por pares:}
\begin{sloppypar}
\begin{itemize}
    \item \textbf{Information Sciences} (Elsevier; IF 6.0) --- 2025. Ciencias de la información, ingeniería del conocimiento, sistemas inteligentes.
    \item \textbf{Information Fusion} (Elsevier; IF 17.4) --- 2023. Fusión de información multisensor y multifuente.
    \item \textbf{IEEE SSCI 2025} --- Simposio sobre Inteligencia Computacional Confiable, Explicable y Responsable (Trondheim, Noruega).
\end{itemize}
\end{sloppypar}
Verificado vía ORCID y confirmación oficial de CMT.

%--- RESEARCH PROJECTS & FUNDING ---
\section*{Proyectos de investigación y financiación}

\begin{itemize}
    \item \textbf{PID2023-150070NB-I00} (MCIN/AEI). \emph{Investigador predoctoral y posdoctoral}. IA Confiable, centrada en explicabilidad e IA segura.
    \item \textbf{Cátedra ENIA IAFER} (EU NextGen, TSI-100927-2023-1). \emph{Investigador predoctoral}. IA Confiable, explicabilidad e IA segura.
    \item \textbf{IAFER CiB --- Cátedra INCIBE de IA Ética para Ciberseguridad} (DaSCI/UGR). \emph{Investigador predoctoral}. XAI aplicada a ciberseguridad e IA ética.
\end{itemize}

%--- TECHNICAL SKILLS ---
\section*{Competencias técnicas}

\blueitem{Lenguajes:} Python, R \\
\blueitem{Entornos:} PyTorch, TensorFlow, scikit-learn \\
\blueitem{Herramientas:} Git, Linux, Conda, TensorBoard, Slurm \\
\blueitem{Computación:} clústeres multi-GPU (NVIDIA A100), entrenamiento distribuido, seguimiento de experimentos

%--- SPOKEN LANGUAGES ---
\section*{Idiomas}

Español --- nativo \quad\textbar\quad Inglés --- fluido

%--- REFERENCES ---
\section*{Referencias}

\begin{minipage}[t]{0.48\textwidth}
    \textbf{Prof.~Francisco Herrera} \\
    \sub{Catedrático de Universidad \\ Dpto.~de CC.~de la Computación e IA, UGR \\ \href{mailto:herrera@decsai.ugr.es}{herrera@decsai.ugr.es}}
\end{minipage}\hfill
\begin{minipage}[t]{0.48\textwidth}
    \textbf{Prof.~Julián Luengo} \\
    \sub{Catedrático de Universidad \\ Dpto.~de CC.~de la Computación e IA, UGR \\ \href{mailto:julianlm@decsai.ugr.es}{julianlm@decsai.ugr.es}}
\end{minipage}

\end{document}
